\documentclass[fancyhdr,bookmark]{ctexbook}
\setCJKmainfont{SimSun}
\setmainfont{Georgia} 	% 設定英文字型
\setromanfont{Georgia} 	% 字型
\setmonofont{Latin Modern Mono}

\usepackage{tikz} % Required for drawing custom shapes
\usepackage[yyyymmdd,hhmmss]{datetime}
\ctexset{today=small}
\usepackage{geometry} 		% 設定邊界
\geometry{
  top=1in,
  inner=1in,
  outer=1in,
  bottom=1in,
  headheight=3ex,
  headsep=2ex
}


\usepackage{xcolor}
\definecolor{ocre}{RGB}{243,102,25} % Define the orange color used for highlighting throughout the book

%\usepackage[x11names,svgnames,dvipsnames]{xcolor}
\usepackage{listings}
\lstset{
	%numbers=left,
	%numberstyle=\tiny,
	basicstyle=\small\ttfamily,
	keywordstyle=\color[rgb]{0.13,0.29,0.53}\textbf,
	commentstyle=\color[rgb]{0.56,0.35,0.01}\textit,
	identifierstyle=\color[rgb]{0.00,0.00,0.00},
	stringstyle=\color[rgb]{0.31,0.60,0.02},
	frame=shadowbox,
	rulesepcolor=\color{red!20!green!20!blue!20},
	backgroundcolor=\color[rgb]{0.97,0.97,0.97},
	tabsize=4,
	breaklines=tr,
	showstringspaces=false,
}
\renewcommand{\lstlistingname}{代码}

%\newcommand{\KeywordTok}[1]{\textcolor[rgb]{0.13,0.29,0.53}{\textbf{{#1}}}}
%\newcommand{\DataTypeTok}[1]{\textcolor[rgb]{0.13,0.29,0.53}{{#1}}}
%\newcommand{\DecValTok}[1]{\textcolor[rgb]{0.00,0.00,0.81}{{#1}}}
%\newcommand{\BaseNTok}[1]{\textcolor[rgb]{0.00,0.00,0.81}{{#1}}}
%\newcommand{\FloatTok}[1]{\textcolor[rgb]{0.00,0.00,0.81}{{#1}}}
%\newcommand{\CharTok}[1]{\textcolor[rgb]{0.31,0.60,0.02}{{#1}}}
%\newcommand{\StringTok}[1]{\textcolor[rgb]{0.31,0.60,0.02}{{#1}}}
%\newcommand{\CommentTok}[1]{\textcolor[rgb]{0.56,0.35,0.01}{\textit{{#1}}}}
%\newcommand{\OtherTok}[1]{\textcolor[rgb]{0.56,0.35,0.01}{{#1}}}
%\newcommand{\AlertTok}[1]{\textcolor[rgb]{0.94,0.16,0.16}{{#1}}}
%\newcommand{\FunctionTok}[1]{\textcolor[rgb]{0.00,0.00,0.00}{{#1}}}
%\newcommand{\RegionMarkerTok}[1]{{#1}}
%\newcommand{\ErrorTok}[1]{\textbf{{#1}}}
%\newcommand{\NormalTok}[1]{{#1}}



\ifxetex
  \usepackage[setpagesize=false, % page size defined by xetex
              unicode=false, % unicode breaks when used with xetex
              xetex]{hyperref}
\else
  \usepackage[unicode=true]{hyperref}
\fi
\hypersetup{breaklinks=true,
            bookmarks=true,
            pdfauthor={An He},
            pdftitle={用Markdown+Pandoc+XeLaTeX写作},
            colorlinks=true,
            urlcolor=blue,
            linkcolor=magenta,
            pdfborder={0 0 0}}
\urlstyle{same}  % don't use monospace font for urls
% Make links footnotes instead of hotlinks:
\renewcommand{\href}[2]{#2\footnote{\url{#1}}}



\title{用Markdown+Pandoc+XeLaTeX写作}
\author{An He}
\date{\today}
\usepackage{cleveref}

\usepackage{fancyhdr}
%\usepackage{lastpage}
\pagestyle{fancy}


\begin{document}
\frontmatter
%----------------------------------------------------------------------------------------
%	TITLE PAGE
%----------------------------------------------------------------------------------------

\begingroup
\thispagestyle{empty}
\begin{tikzpicture}[remember picture,overlay]
\node[inner sep=0pt] (background) at (current page.center) {\includegraphics[width=\paperwidth]{Pictures/background}};
\draw (current page.center) node [fill=ocre!30!white,fill opacity=0.6,text opacity=1,inner sep=1cm]
{\Huge\centering\bfseries\sffamily\parbox[c][][t]{\paperwidth}
{\centering 用Markdown+Pandoc+XeLaTeX写作\\[13pt] % Book title
{\huge An He} % Author name
}};
\end{tikzpicture}
\vfill
\endgroup
\newpage
\newpage


%----------------------------------------------------------------------------------------
%	COPYRIGHT PAGE
%----------------------------------------------------------------------------------------
\newpage
~\vfill
\thispagestyle{empty}

\noindent Copyright \copyright\ \the\year\  An He\\ % Copyright notice

\noindent \textsc{Published by \LaTeX}\\ % Publisher
\noindent \textsc{https://github.com/annProg/pandoc-template}\\ % URL

\noindent Licensed under the Creative Commons Attribution-NonCommercial 3.0 Unported License (the ``License''). You may not use this file except in compliance with the License. You may obtain a copy of the License at \url{http://creativecommons.org/licenses/by-nc/3.0}. Unless required by applicable law or agreed to in writing, software distributed under the License is distributed on an \textsc{``as is'' basis, without warranties or conditions of any kind}, either express or implied. See the License for the specific language governing permissions and limitations under the License.\\ % License information

\noindent \textit{最后编译日期, \today\ \currenttime } % Printing/edition date


    
\chapter*{序言}
\addcontentsline{toc}{chapter}{序言}
使用Markdown+Pandoc生成LaTeX文档


{
\hypersetup{linkcolor=black}
\setcounter{tocdepth}{2}
\tableofcontents
}



\mainmatter
    % 在此命令之后的页码为阿拉伯数字
    % 以下为正文
\chapter{代码块}\label{ux4ee3ux7801ux5757}

pandoc扩展的Markdown语法可以为代码块添加ID属性及语言类型属性，形如\lstinline!{#id .language}!，其中ID属性可以用来做交叉引用，使用\lstinline!\ref{id}!在正文中引用代码块。例如代码块\ref{code:demo}。

\begin{lstlisting}[language=bash, label=code:demo, caption=code:demo, float=htbp]
~~~{#code:demo .bash}
#!/bin/bash
echo "Hello LaTeX"
echo "Hello Pandoc"
echo "Hello Markdown"
~~~~
\end{lstlisting}

\chapter{图片}\label{ux56feux7247}

\backmatter

\end{document}
